%=====================================================================
\chapter{The Checkout and the Order State Machine}\label{ch:checkout}
%=====================================================================
\locator{3}{Part III --- From Basket to Order}

\begin{outcomes}
By the end of this chapter you will be able to:
\begin{tight}
  \item \textbf{Name} the steps of a checkout in order, and \textbf{state}
        what each one validates and what it saves.
  \item \textbf{Explain} why an order is a state machine rather than a record,
        and \textbf{read} a shop's order statuses as the states of one.
  \item \textbf{Measure} the cost of each checkout step in time, queries and
        writes, and \textbf{identify} where an order is actually created.
  \item \textbf{Diagnose} a checkout that fails silently, using the shop's own
        responses rather than its screens.
  \item \textbf{Judge} where a checkout can safely be shortened, and where
        shortening it would lose something the shop needs.
\end{tight}
\end{outcomes}

\begin{prereq}
\chref{basket}: where the basket lives, and that it is rows keyed by a session
identifier. \chref{catalogue}: reading the queries a page actually ran.
Chapter~5's container, with a shop you can place a test order in.
\end{prereq}

\begin{keyterms}
checkout $\bullet$ checkout step $\bullet$ guest checkout $\bullet$ order
$\bullet$ order status $\bullet$ state machine $\bullet$ order history
$\bullet$ confirmation $\bullet$ idempotency $\bullet$ abandonment $\bullet$
conversion funnel $\bullet$ write amplification $\bullet$ silent failure
\end{keyterms}

\section{A Checkout Is a Sequence, Not a Page}

The basket in \chref{basket} is a list of rows. An \term{order} is something
else entirely: a commitment, by a named person, to pay a stated amount for
stated goods delivered to a stated address. Turning the first into the second
is what a checkout does, and it cannot be done in one step because each piece
of information constrains the next --- the delivery address decides which
shipping methods exist, the shipping method decides the total, and the total
decides what the customer is agreeing to.

\begin{definitionbox}[Order, Order Status, Order History]
An \term{order} is a row recording an intention to buy, together with rows for
the products, the totals and the addresses as they were \emph{at the moment of
ordering}. It does not point at the catalogue; it copies from it, because the
price of a product next month must not change what a customer agreed to pay
today.

\smallskip
An \term{order status} is the state the order is in: pending, processing,
shipped, complete, refunded. The set of statuses is the shop's, not the
standard's.

\smallskip
The \term{order history} is the list of status changes, each with a date and
an optional comment. The status column says where the order is; the history
says how it got there, and it is the only part a dispute can be settled from.
\end{definitionbox}

\dsfig{%
\begin{tikzpicture}[
  stp/.style={rounded corners=2pt, draw=#1, line width=1.0pt, fill=#1!8,
              text=#1!60!black, font=\bfseries\scriptsize, align=center,
              text width=2.25cm, minimum height=0.85cm},
  arr/.style={-{Stealth[length=2.0mm]}, primary!55, line width=0.9pt},
  tag/.style={font=\scriptsize, text=neutral, align=center, text width=2.6cm},
  band/.style={rounded corners=3pt, draw=#1!40, line width=0.8pt,
               fill=#1!4, inner sep=0pt}
]
  \node[stp=secondary] (b)  at (0,0)    {Basket};
  \node[stp=secondary] (r)  at (2.6,0)  {Who you are};
  \node[stp=secondary] (sh) at (5.2,0)  {Where to};
  \node[stp=secondary] (sm) at (7.8,0)  {How to send};
  \node[stp=secondary] (pm) at (10.4,0) {How to pay};
  \node[stp=accent]    (cf) at (2.6,-2.1) {Confirm};
  \node[stp=accent]    (py) at (5.2,-2.1) {Pay};
  \node[stp=greenhl]   (sc) at (7.8,-2.1) {Success};

  \foreach \a/\b in {b/r, r/sh, sh/sm, sm/pm}{\draw[arr] (\a) -- (\b);}
  \draw[arr] (pm.south) -- ++(0,-0.75) -| (cf.north);
  \draw[arr] (cf) -- (py);
  \draw[arr] (py) -- (sc);

  \node[tag, text=secondary] at (5.2,0.95)
       {\textbf{Typing.} Nothing is committed. The customer may leave and
        nothing is lost.};
  \node[tag, text=accent, text width=3.4cm] at (3.9,-3.15)
       {\textbf{The order row is written here}, at status~0, \emph{before}
        payment.};
  \node[tag, text=greenhl, text width=2.9cm] at (7.8,-3.1)
       {Status becomes Pending. Only now has anything happened.};
  \node[font=\scriptsize\itshape, text=neutral, align=left, text width=4.0cm]
       at (10.6,-2.4)
       {A customer who stops between Confirm and Pay leaves a row behind that
        no screen will show you.};
\end{tikzpicture}}%
{The guest checkout as the shop actually runs it. The first five steps collect;
the last three commit. The gap between the order row appearing and the payment
confirming it is where orders go missing.}{checkoutsteps}

\section{The Fault That Stops Every Order}

Before any of that could be measured, the checkout had to be made to work, and
on a freshly installed shop it does not. This is worth the space because the
way the failure presents itself is the lesson, not the fault itself.

\begin{alertbox}[A stock install cannot complete a guest checkout]
Driving the sequence by hand against the shop built in \chref{stack}, every
step succeeded until the shipping quote, which answered:

\smallskip
\hspace*{1em}\texttt{\{"error":"Shipping address required!"\}}

\smallskip
--- although a shipping address had been supplied and accepted one step
earlier, with \texttt{\{"success":"\dots\ has been saved!"\}}.

\smallskip
\textbf{The cause.} The guest registration step copies the billing address to
the shipping address, but only inside a branch guarded by the setting
\texttt{config\_checkout\_payment\_address}. That setting is read by
\textbf{ten} controllers in \texttt{catalog/}, and appears as a checkbox in the
admin settings form --- but \textbf{the installer never writes it}. An absent
setting reads as false, the branch never runs, no shipping address is ever
stored, and the quote has nothing to work from.

\smallskip
\textbf{The fix} is to open \adminpath{System, Settings, Edit, Option}, tick
\uiel{Payment Address required} and save, which writes every field on that
form, including the missing one. One checkbox.

\smallskip
\textbf{What makes this worth a page.} No page returned an error code. No log
recorded a failure. Every screen rendered. A shopkeeper would see customers
reaching the shipping step and leaving, and would reasonably conclude that the
shipping charges were too high. This is the same shape of fault as the store
identifier in \chref{catalogue}: \textbf{the system did not break, it answered
a different question}, and only driving it end to end and reading the replies
found it.
\end{alertbox}

\begin{conceptbox}[How to find a fault of this shape]
\begin{tightnum}
  \item \textbf{Read the reply, not the page.} Each step answers with JSON. The
        error was in the reply all along; the browser merely rendered it in a
        place nobody looks.
  \item \textbf{Trust the sequence, not the screens.} Drive the steps in order
        from a script. The step that first disagrees with the previous one is
        the step to read.
  \item \textbf{Then read the source for that step only.} Not before --- the
        measurement tells you which of four thousand files to open.
  \item \textbf{Change one thing and repeat.} Setting the one flag made the
        whole sequence complete, which is what turns a hypothesis into a cause.
\end{tightnum}
\end{conceptbox}

\section{What Each Step Costs}

With the checkout working, the sequence can be measured. The question put to it
is a specific one: \emph{is the last step the expensive one?} It is the common
belief, and it is worth settling with numbers.

\begin{measured}[A complete guest checkout, step by step]
Produced by \texttt{code/m10\_checkout.py --repeat 5} against the stack in
Appendix~A: medians of five complete checkouts, each buying two products.
Queries and writes are counted from MariaDB's slow query log, not estimated.

\rowcolors{2}{white}{lightbg}
\begin{longtable}{@{}L{3.4cm}L{1.7cm}L{1.9cm}L{1.8cm}L{1.7cm}L{1.6cm}@{}}
\toprule
\rowcolor{primary!15}
\textbf{Step} & \textbf{Time} & \textbf{Queries} & \textbf{Writes} &
\textbf{DB time} & \textbf{Share} \\
\midrule
\endhead
fill basket & 55\,ms & 57 & 6 & 23\,ms & 11\% \\
register (guest) & 30\,ms & 41 & 3 & 9\,ms & 6\% \\
shipping address$^*$ & 21\,ms & 32 & 2 & 6\,ms & 4\% \\
shipping method & 56\,ms & 67 & 4 & 22\,ms & 11\% \\
payment method & 54\,ms & 74 & 4 & 17\,ms & 10\% \\
\textbf{confirm} & \textbf{149\,ms} & \textbf{125} & \textbf{25} & 93\,ms &
\textbf{28\%} \\
pay & 64\,ms & 102 & 8 & 25\,ms & 12\% \\
success & 94\,ms & 109 & 3 & 37\,ms & 18\% \\
\midrule
\textbf{Total} & \textbf{523\,ms} & \textbf{607} & \textbf{55} & --- & \\
\bottomrule
\end{longtable}

\smallskip
$^*$ \emph{A guest never uses this step.} It redirects to the login page,
because its guard requires a logged-in customer. The round trip costs 21\,ms
and 32 queries and achieves nothing; a guest's addresses were both set by the
registration step.

\smallskip
\textbf{The common belief is wrong as stated.} No single step costs more than
the rest together. \texttt{confirm} is the most expensive at 28\%, not 50\%.

\smallskip
\textbf{But split the sequence where the customer stops typing:}

\smallskip
\begin{tabular}{@{}L{5.2cm}L{2.6cm}L{2.6cm}L{2.6cm}@{}}
& \textbf{Time} & \textbf{Queries} & \textbf{Writes} \\
\textbf{Typing} (5 steps) & 217\,ms & 271 & 19 \\
\textbf{Ordering} (3 steps) & 307\,ms & 336 & \textbf{36} \\
\end{tabular}

\smallskip
\textbf{The three steps after the last keystroke cost more than the five before
it} --- 40\% more time, and nearly twice the writes. That is the defensible
version of the claim.

\smallskip
\textbf{The sharpest number in the table is the writes column.}
\texttt{confirm} performs \textbf{25 of the 55 writes} in the entire checkout,
against two to six for every other step. It is not slow because it is
complicated; it is slow because it is the step that \emph{writes the order}:
the order row, the product rows, the totals, the addresses, all copied out of
the basket and the session so that they can never change again.

\smallskip
\textbf{For scale:} one checkout is \textbf{607 queries}. A category page in
\chref{catalogue} was 144. A single customer buying two items costs the
database more than four people browsing.
\end{measured}

\section{An Order Is a State Machine}

The reason \texttt{confirm} and \texttt{pay} are separate steps --- and the
reason the gap between them matters --- is that an order has a life after it is
created, and that life is a sequence of states.

\begin{examplebox}[The fourteen states this shop defines]
\adminpath{System, Localisation, Order Statuses} lists them. They are not a
standard; they are this shop's vocabulary, and a shop may add its own.

\smallskip
\begin{tabular}{@{}L{4.4cm}L{4.4cm}L{4.4cm}@{}}
1 Pending & 8 Denied & 13 Chargeback \\
2 Processing & 9 Canceled Reversal & 14 Expired \\
3 Shipped & 10 Failed & 15 Processed \\
5 Complete & 11 Refunded & 16 Voided \\
7 Canceled & 12 Reversed & \\
\end{tabular}

\smallskip
Three things are worth noticing. \textbf{There is no 4 and no 6} --- the
numbers are identifiers, not an order. \textbf{Several are about money moving
backwards} (Refunded, Reversed, Chargeback, Canceled Reversal), which is
\chref{orders}'s subject and a reminder that selling is the easy half.
\textbf{And there is no zero.}
\end{examplebox}

\begin{measured}[Status 0, and the orders nobody can see]
From \texttt{code/m10\_checkout.py --abandon}, which runs the checkout as far
as \texttt{confirm} and then stops, exactly as a customer who changes their
mind at the payment button does.

\smallskip
\begin{tabular}{@{}L{6.6cm}L{6.6cm}@{}}
Orders at status 0 before & 0 \\
Orders at status 0 after & \textbf{1} \\
Order row written? & \textbf{yes}, total \$854.40 \\
Visible in the admin order list? & \textbf{no} \\
\end{tabular}

\smallskip
\textbf{The order row is written at \texttt{confirm}, before any payment, and
left at status 0.} Status~0 is not one of the fourteen. The admin order list
filters it out. Nothing will ever remove it.

\smallskip
This is not a bug; it is a deliberate design, and the reason for it is sound.
The order number has to exist before the customer is sent to a payment
provider, because the provider must be told what it is being paid for, and the
answer that comes back has to find its way to something. The cost of that
design is a table containing rows for orders that never happened.

\smallskip
\textbf{Why it matters to you:} a shop owner reporting ``we have orders in the
database that aren't in the orders list'' has not found a fault. They have
found their abandonment rate. \chref{dashboard} turns that into a number worth
watching, and the ratio of status-0 rows to real orders is one of the most
useful figures a small shop can compute from data it already has.
\end{measured}

\begin{alertbox}[Two requests, one order --- the problem confirm cannot solve
alone]
If the customer's connection drops after \texttt{confirm} is sent but before
the reply arrives, their browser may retry. If the shop treats the second
request as a new one, the customer is charged twice.

\smallskip
The property that prevents this is \term{idempotency}: performing an operation
twice has the same effect as performing it once. OpenCart approaches it by
keeping the order identifier in the session and updating that order rather than
creating another, which is why \texttt{confirm} can be called repeatedly in a
session without multiplying orders --- and why the order identifier must be
established \emph{before} payment, not after.

\smallskip
It is not a complete answer, and \chref{gateway} returns to it where it is
hardest: when the second request comes not from the customer but from the
payment provider, minutes later, over a connection that is not the customer's
at all.
\end{alertbox}

\section{Where a Checkout Can Be Shortened, and Where It Cannot}

Every step is a place to lose a customer, so the pressure to remove steps is
constant. The measurements say which removals are free.

\begin{conceptbox}[Four changes, ranked by what the evidence supports]
\begin{tightnum}
  \item \textbf{Remove the guest's shipping-address round trip.} It is measured
        at 21\,ms and 32 queries, it redirects to a login page the guest cannot
        use, and the data it would set is already set. Pure waste, and the only
        change here with no trade-off.
  \item \textbf{Merge the shipping and payment method steps.} 110\,ms and 141
        queries between them, and both are selections from short lists. The
        cost is that a shipping choice can change the total that the payment
        step displays, so they must then be recomputed together.
  \item \textbf{Skip the address entirely for digital goods.} There is nothing
        to deliver. This is a real saving and the platform supports it ---
        \texttt{hasShipping()} is already consulted.
  \item \textbf{Remove the confirmation screen.} Tempting --- 28\% of the time.
        \textbf{Do not.} It is the last point at which the customer sees what
        they are agreeing to pay, and removing it moves every mistake from
        before the sale to after it, where it costs a refund, a complaint and
        \chref{legal}'s attention rather than a click.
\end{tightnum}
\end{conceptbox}

\begin{pitfall}
\begin{tight}
  \item \textbf{Testing a checkout only through the browser.} The fault at the
        start of this chapter was invisible there and obvious in the replies.
  \item \textbf{Assuming the slow step is the last step.} It is the most
        expensive single step, at 28\%, which is not the same claim.
  \item \textbf{Counting orders without filtering status 0,} and reporting
        revenue that includes abandoned baskets.
  \item \textbf{Deleting status-0 rows to tidy up.} They are the audit trail of
        every customer who reached the payment button, and a payment provider's
        late reply may still be looking for one of them.
  \item \textbf{Pointing an order's product rows at the catalogue} instead of
        copying names and prices into the order. The catalogue changes; what a
        customer agreed to must not.
  \item \textbf{Removing the confirmation screen to raise conversion.} It
        raises conversion and refunds together.
\end{tight}
\end{pitfall}

\begin{lab}[Drive a checkout, then break it on purpose]
\begin{lstlisting}[language=bash]
# --- 1. Make the shop able to take an order at all -------------------
# Admin: System > Settings > Edit > Option. Tick "Payment Address
# required" and save. Write down what you expect that to change.

# --- 2. Place one order by hand, in the browser ----------------------
# Note how many pages you visited and roughly how long it took.

# --- 3. Now measure it --------------------------------------------
cd code
python m10_checkout.py --repeat 5

# --- 4. Predict before you look -------------------------------------
# Which step has the most WRITES? Most people say "success".

# --- 5. Watch an order appear before it is paid for -----------------
python m10_checkout.py --abandon

# --- 6. Look for it in the admin panel ------------------------------
# Admin: Sales > Orders. It is not there. Now look in the table:
docker compose exec db mariadb -ushop -pshoppw opencart \
  -e "SELECT order_id, order_status_id, total FROM oc_order;"

# --- 7. Read the state machine off your own data --------------------
python m10_checkout.py --states

# --- 8. Move an order by hand through three states ------------------
# Admin: Sales > Orders > View > History. Pending -> Processing ->
# Shipped. Then read oc_order_history and say what the status column
# alone would not have told you.

# --- 9. Clean up ----------------------------------------------------
python m10_checkout.py --clean
\end{lstlisting}

\smallskip
\textbf{What to notice.} Step 1 is the whole of \chref{planning}'s argument in
one checkbox: a platform you did not write will have defects, and the skill
that matters is finding them from the outside, not avoiding platforms.
\end{lab}

\begin{checkpoint}
\begin{tightnum}
  \item A shop reports that customers abandon at the shipping step. Describe,
        in order, the three things you would do before accepting that the
        shipping charges are too high --- and say what the first chapter
        measurement would have told you.
  \item \texttt{confirm} is 28\% of the checkout's time but 45\% of its writes.
        Explain why those two shares differ, and what that implies for which
        step you would optimise if the shop became ten times busier.
  \item Your client asks you to delete every order at status 0 because they
        ``clutter the database''. Give two distinct reasons to refuse, and one
        thing you would build instead.
\end{tightnum}
\end{checkpoint}

\begin{chaptersummary}
\begin{tight}
  \item A checkout is a sequence because each answer constrains the next. An
        order copies what it needs out of the catalogue rather than pointing at
        it, so that later price changes cannot rewrite an agreement.
  \item Measured: a stock install \emph{cannot complete a guest checkout}. The
        setting \texttt{config\_checkout\_payment\_address} is read by ten
        controllers and written by none, so no shipping address is ever stored.
        Every page still returned 200.
  \item The way to find a fault of that shape is to drive the steps from a
        script and read the replies, then open the source for the one step that
        disagreed.
  \item Measured: 523\,ms and 607 queries for one complete checkout --- more
        than four category pages.
  \item Measured: no single step costs more than the rest together, but the
        three steps after the last keystroke cost 307\,ms against 217\,ms and
        nearly twice the writes.
  \item Measured: \texttt{confirm} performs 25 of the checkout's 55 writes,
        because it is the step that creates the order.
  \item Measured: abandoning at the payment button leaves an order row at
        status 0 --- a status the shop does not define, that the admin list
        hides, and that nothing removes. It is the abandonment rate, not a
        fault.
  \item The order history, not the status column, is what a dispute is settled
        from.
  \item A step may be removed when the evidence says it does nothing. The
        confirmation screen is not that step.
\end{tight}
\end{chaptersummary}

\begin{reviewq}
  \item \textbf{(MCQ)} In the measured checkout, which step performs the most
        database writes? (a)~fill basket; (b)~payment method; (c)~confirm;
        (d)~success.
  \item \textbf{(MCQ)} An order sitting at status 0 means: (a)~the order
        failed; (b)~the order was refunded; (c)~the order row was created but
        payment was never confirmed; (d)~the status table is corrupt.
  \item \textbf{(Short)} Explain why an order copies the product name and price
        instead of referring to the catalogue row, and give one thing that
        would go wrong if it did not.
  \item \textbf{(Short)} Define idempotency, and explain why a checkout needs
        it in terms of a specific thing that would otherwise happen to a
        customer.
  \item \textbf{(Short)} The stock shop answered every request with 200 while
        being unable to take a single order. Explain what this says about
        monitoring a shop by availability.
  \item \textbf{(Applied)} You are given a shop with 1{,}200 rows in
        \texttt{oc\_order}, of which 820 are at status 0. Say what you can and
        cannot conclude from that ratio, what two further figures you would
        want, and what you would recommend --- with the reasoning for each
        step shown.
  \item \textbf{(Applied)} Design a measurement that would tell you at which
        checkout step customers actually leave. State what you would record,
        where in the sequence you would record it, how you would avoid counting
        the same visitor twice, and how long you would need to run it before
        the answer meant anything for a shop taking thirty orders a day.
\end{reviewq}
